\section{Proposed Method}
In this section, we present the architecture and theoretical formulation of the Cross-Modality Semantic Alignment Transformer (CMSA-Trans). The core objective of the proposed framework is to bridge the gap between high-dimensional vibration signals and abstract technical concepts by using the extensive mechanical knowledge embedded in Large Language Models (LLMs). By mapping physical signal features into a shared semantic space, the model enables zero-shot fault diagnosis of rotating machinery, where the health states of unseen components can be recognized without prior training samples.

\subsection{Overall Architecture}
The overall architecture of the proposed CMSA-Trans is illustrated in Fig.~\ref{fig:framework}. The CMSA-Trans framework is designed to transform raw sensor data into semantic-aligned representations that are directly comparable to linguistic descriptions of machinery faults. The high-level workflow begins with the acquisition of raw vibration signals from sensors mounted on rotating components, such as bearings or gearboxes. These signals are typically non-stationary and contain periodic impulsive components that signify specific fault modes.

The architecture consists of four primary functional modules: a Signal Encoder, a Semantic Prototyping module, a Cross-Modality Alignment layer, and a Zero-Shot Classification head. The Signal Encoder uses a transformer-based structure \cite{vaswani2017attention, zerveas2021transformer} to extract deep temporal features from the vibration data. Simultaneously, the Semantic Prototyping module interfaces with a pre-trained LLM \cite{radford2021learning} to generate high-fidelity technical descriptions for each fault category, which are then converted into fixed-dimensional semantic vectors (prototypes). This automation removes the need for expert-defined attribute matrices.

The heart of the CMSA-Trans is the Cross-Modality Alignment module, which employs a mutual-attention mechanism to project the learned signal features into the semantic space. Unlike traditional diagnostic models that map signals directly to discrete class labels, our approach maps signals to the continuous semantic space defined by the LLM. This allows the model to "reason" about the similarity between a physical signal pattern and a technical concept (e.g., "outer race pitting"). This alignment process is inspired by recent advances in contrastive learning \cite{radford2021learning, oord2018representation, chen2020simple}. Finally, the zero-shot classification is performed by identifying the nearest semantic prototype for a given signal embedding. This structured approach ensures that the model captures the underlying physics of machine failures, enabling generalization to new, unseen fault modes.

\subsection{Signal Encoder: Time-Series Transformer}
Vibration signals from rotating machinery are characterized by their multi-scale periodicity and non-stationary impulses. Traditional deep learning backbones, such as Convolutional Neural Networks (CNNs), often fail to capture the long-range dependencies between impacts occurring across multiple shaft revolutions due to their local receptive fields. While Recurrent Neural Networks (RNNs) can handle sequential data, they suffer from vanishing gradients when processing long industrial signal sequences. To address these limitations, we adopt a Time-Series Transformer (TST) as the Signal Encoder \cite{zerveas2021transformer}, which uses self-attention to maintain a global view of the entire signal window.

Given a raw vibration signal segment $\mathbf{X} \in \mathbb{R}^{L \times 1}$, where $L$ is the window length, we first partition the signal into $N$ non-overlapping patches $\mathbf{x}_p \in \mathbb{R}^{P \times 1}$, where $P$ is the patch size. Each patch represents a localized temporal window of the vibration waveform. These patches are then projected into a high-dimensional embedding space using a linear layer with weights $\mathbf{W}_b \in \mathbb{R}^{P \times d_m}$, where $d_m$ is the model dimension. This projection allows the model to process 1D temporal segments as discrete tokens. To preserve the relative and absolute temporal order of the signal, learnable positional encodings $\mathbf{E}_{pos} \in \mathbb{R}^{N \times d_m}$ are added to the projected embeddings. The resulting sequence of input embeddings $\mathbf{Z}_0$ is formulated as follows:
\begin{equation}
\mathbf{Z}_0 = [\mathbf{x}_p^{(1)}\mathbf{W}_b; \mathbf{x}_p^{(2)}\mathbf{W}_b; \dots; \mathbf{x}_p^{(N)}\mathbf{W}_b] + \mathbf{E}_{pos}
\end{equation}
where $\mathbf{Z}_0 \in \mathbb{R}^{N \times d_m}$ serves as the input to the transformer backbone.

The embedded sequence passes through a stack of transformer encoder layers. Each layer consists of a Multi-Head Self-Attention (MHSA) sub-layer and a Position-wise Feed-Forward Network (FFN). The MHSA mechanism is particularly effective for vibration signals because it allows each segment of the signal to attend to every other segment, regardless of their distance in time. This is critical for detecting fault signatures like inner-race faults, where the impulsive pattern is modulated by the shaft rotation and may be separated by long intervals of background noise. The FFN layer then applies non-linear transformations to each token independently, refining the feature representation. Residual connections and layer normalization are employed to ensure stable training. The output of the final transformer layer, denoted as $\mathbf{H}_v \in \mathbb{R}^{N \times d_m}$, represents the refined long-range structural features. This encoder effectively mitigates the inductive bias of local-receptive field models and enables more robust feature extraction in the presence of industrial background noise. Through this architecture, the TST captures the "rhythm" of the machine, which is fundamental for subsequent semantic alignment.

\subsection{Semantic Prototyping via LLM}
A major bottleneck in zero-shot fault diagnosis is the construction of a discriminative semantic space. Conventional methods rely on manual attribute engineering (e.g., binary labels for "impulsive" or "harmonic"), which is labor-intensive, subjective, and fails to scale to complex industrial systems with hundreds of failure modes. Our CMSA-Trans overcomes this by automating semantic generation through an LLM interface, treating the LLM as a virtual domain expert.

We use the world knowledge of a pre-trained LLM (e.g., GPT-4) to generate detailed technical descriptions for each fault mode. Specifically, for each seen class in the training set and unseen class in the test set, we provide a structured prompt such as: "Provide a detailed technical description of the vibration characteristics for a [Fault Name] in a bearing, focusing on its frequency components, impact patterns, and resonance bands." The LLM generates a text corpus that describes the physical manifestations of the fault, effectively linking the label to mechanical engineering principles. This process is inherently superior to manual labeling because the LLM can synthesize information from a vast array of technical manuals and research papers.

Table~\ref{tab:llm_semantics} presents a comparison of the LLM-generated semantic descriptions with traditional word-level embeddings. These textual descriptions are then transformed into fixed-length semantic vectors (prototypes) using a sentence-level embedding model such as SBERT \cite{reimers2019sentence}. For a class $y$, the LLM-generated text $T_y$ is encoded into a semantic prototype $\mathbf{a}_y \in \mathbb{R}^{d_s}$, where $d_s$ is the dimensionality of the semantic space. These prototypes $\mathbf{a}_y$ act as anchors in the embedding space. Since the LLM uses external technical knowledge to define these anchors, the resulting semantic space is inherently grounded in the physics of failures. This allows the model to distinguish between similar fault types based on subtle differences in their technical descriptions, such as the specific Ball Pass Frequency (BPF) characteristics or sideband patterns mentioned in the text. Also, because these prototypes are generated in a continuous vector space, the model can represent the relationships between faults—for instance, two different seal-rubbing faults will naturally have similar semantic vectors. This rich semantic structure is what enables the cross-modal alignment module to generalize signal features to previously unobserved conditions.

\subsection{Cross-Modality Alignment}
The core innovation of the proposed method is the Cross-Modality Alignment module, which maps the vibration signal features $\mathbf{H}_v$ into the semantic space $\mathcal{S}$ populated by the LLM-derived prototypes through a dynamic attention-based process.

We use a mutual-attention mechanism where the signal features $\mathbf{H}_v$ act as the query, and the semantic prototypes of the seen classes $\mathbf{A}_{seen} \in \mathbb{R}^{C_{seen} \times d_s}$ act as the keys and values. This allows the model to learn cross-modal sensitivity: the attention mechanism assigns higher weights to signal tokens exhibiting characteristics described in the LLM text (e.g., "high-frequency impact"). We define the Query ($\mathbf{Q}$), Key ($\mathbf{K}$), and Value ($\mathbf{V}$) matrices for the alignment layer as:
\begin{equation}
\mathbf{Q} = \text{Norm}(\mathbf{H}_v) \mathbf{W}_Q, \quad \mathbf{K} = \mathbf{A}_{seen} \mathbf{W}_K, \quad \mathbf{V} = \mathbf{A}_{seen} \mathbf{W}_V
\end{equation}
where $\mathbf{W}_Q$, $\mathbf{W}_K$, and $\mathbf{W}_V$ are learnable weight matrices. The cross-modality attention score is then computed to aggregate the signal features aligned with the semantic context:
\begin{equation}
\mathbf{H}_{aligned} = \text{softmax}\left(\frac{\mathbf{Q}\mathbf{K}^T}{\sqrt{d_k}}\right) \mathbf{V}
\end{equation}
where $\mathbf{H}_{aligned}$ represents the semantic-aware signal representation and $d_k$ is the scaling factor. This ensures that the alignment is driven by the intrinsic relationship between signal morphology and technical terminology.

Following the attention layer, we apply global average pooling to the aligned sequence and pass it through a final projection layer $\mathbf{W}_p \in \mathbb{R}^{d_m \times d_s}$. The final signal embedding $\mathbf{v}_x \in \mathbb{R}^{d_s}$ in the semantic space is given by:
\begin{equation}
\mathbf{v}_x = \text{ReLU}(\text{Pooling}(\mathbf{H}_{aligned}) \mathbf{W}_p)
\end{equation}
By mapping $\mathbf{v}_x$ into the same $d_s$-dimensional space as the LLM prototypes, we create a unified platform for direct similarity measurement between signals and textual concepts.

\subsection{Training and Optimization}
The training phase aims to align the signal embedding $\mathbf{v}_x$ with its corresponding semantic prototype $\mathbf{a}_{y}$ for seen classes. We employ a contrastive alignment loss \cite{chen2020simple} based on cosine similarity to pull the signal embedding towards its correct prototype while pushing it away from all other prototypes in the seen set. This enables the learning of an embedding space where clusters are formed based on health-state semantics. This contrastive approach is robust against intra-class variations, as it focuses on the invariant semantic features.

For a training batch of size $B$, the total loss function $\mathcal{L}_{total}$ is defined as a multi-class cross-entropy variant applied to cosine distances:
\begin{equation}
\mathcal{L}_{total} = -\frac{1}{B} \sum_{i=1}^{B} \log \frac{\exp(\text{cos}(\text{v}_x^{(i)}, \mathbf{a}_{y_i}) / \tau)}{\sum_{j \in \mathcal{Y}_{seen}} \exp(\text{cos}(\mathbf{v}_x^{(i)}, \mathbf{a}_j) / \tau)}
\end{equation}
where $\text{cos}(\cdot)$ is the cosine similarity between two vectors, $\tau$ is a temperature hyperparameter, and $\mathcal{Y}_{seen}$ is the set of seen class labels. The temperature parameter controls the sharpness of the probability distribution over classes.

During inference, a test vibration signal $\mathbf{X}_{test}$ is passed through the signal encoder and alignment module to obtain its semantic embedding $\mathbf{v}_{test}$. The zero-shot classification is then performed by identifying the nearest semantic prototype among all candidate classes, including unseen ones. The predicted label $\hat{y}$ is determined by the following utility function:
\begin{equation}
\hat{y} = \arg\max_{c \in \mathcal{Y}_{all}} \text{cos}(\mathbf{v}_{test}, \mathbf{a}_c)
\end{equation}
where $\mathcal{Y}_{all} = \mathcal{Y}_{seen} \cup \mathcal{Y}_{unseen}$ encompasses both seen and unseen fault categories, and $\mathbf{a}_c$ denotes the LLM-generated semantic prototype for class $c$. This formulation enables the model to generalize to new fault types without retraining, as the semantic prototypes for unseen classes are generated independently by the LLM.

